% GENERATED FILE. Edit canonical lessons, then run npm run generate:books.
% canonical-source-hash: fnv1a64:cd37786c371b7900
% canonical-lessons: GU-C388-ruprekha, GU-C388-upakram, GU-C388-abhiyan, GU-C388-arambhvu, GU-C388-adarvu

\chapter{Getting Started}
\label{ch:gu-getting-started}
\hlchaptermodality{388}

\begin{chapteropening}
\textbf{By the end of this chapter:} \emph{I can say what we will begin and what it will take.}

Complete the last lesson of chapter 388: I can say what we will begin and what it will take.
\end{chapteropening}

\section[\texorpdfstring{\gu{રૂપરેખા} (rūprekhā)}{રૂપરેખા (rūprekhā)}]{\gu{રૂપરેખા} (rūprekhā) — an outline}

\label{lesson:GU-C388-ruprekha}



\begin{quote}
Before the new one: say the Gujarati for lovely, then the Gujarati for to deteriorate.
\end{quote}

\subsection*{You'll want to know: \gu{રૂપરેખા}}
\textbf{\gu{રૂપરેખા}} — \emph{rūprekhā} — \textquotedblleft{}an outline\textquotedblright{}.

\begin{quote}
\gu{જો} \gu{રૂપરેખા} \gu{તૈયાર} \gu{હશે} \gu{તો} \gu{કામ} \gu{ઝડપી} \gu{થશે}.
\end{quote}

\begin{quote}
\emph{jo rūprekhā taiyār haśe to kām jhaḍpī thaśe.}
\end{quote}

\begin{quote}
If the outline is ready, the work will go faster.
\end{quote}

A feminine noun.

\begin{cousinweb}
\textbf{\gu{રૂપરેખા}} (\emph{rūprekhā}) is built from Sanskrit \emph{rūpa}, \textquotedblleft{}form, shape\textquotedblright{}, and \emph{rekhā}, \textquotedblleft{}a line\textquotedblright{} (Monier-Williams): the lines of a shape.
\end{cousinweb}

\begin{grammarlens}[title={What you've built}]
One more word for saying what happens if.
\end{grammarlens}

\subsection*{Your turn}
\begin{itemize}
\raggedright
  \item \emph{Say it:} \emph{rūprekhā}
  \item \emph{Say it:} \emph{rūprekhā}, once more
  \item \emph{Say it:} \emph{jo rūprekhā taiyār haśe to kām jhaḍpī thaśe.} — If the outline is ready, the work will go faster.
  \item \emph{Recall:} say the Gujarati for lovely, then the Gujarati for to deteriorate, then say \emph{rūprekhā} again
\end{itemize}

\begin{tcolorbox}[breakable,colback=teal!4,colframe=teal!35!black,title={Before you move on}]
What does \gu{રૂપરેખા} mean? (\textquotedblleft{}An outline\textquotedblright{}.) Say it once more.
\end{tcolorbox}
\section[\texorpdfstring{\gu{ઉપક્રમ} (upakram)}{ઉપક્રમ (upakram)}]{\gu{ઉપક્રમ} (upakram) — an initiative}

\label{lesson:GU-C388-upakram}



\begin{quote}
Before the new one: say the Gujarati for to deteriorate, then the Gujarati for an outline.
\end{quote}

\subsection*{You'll want to know: \gu{ઉપક્રમ}}
\textbf{\gu{ઉપક્રમ}} — \emph{upakram} — \textquotedblleft{}an initiative\textquotedblright{}.

\begin{quote}
\gu{જો} \gu{આ} \gu{ઉપક્રમ} \gu{સફળ} \gu{થશે} \gu{તો} \gu{બીજાં} \gu{ગામો} \gu{પણ} \gu{તે} \gu{અપનાવશે}.
\end{quote}

\begin{quote}
\emph{jo ā upakram saphaḷ thaśe to bījā̃ gāmo paṇ te apnāvśe.}
\end{quote}

\begin{quote}
If this initiative succeeds, other villages will take it up too.
\end{quote}

A masculine noun.

\begin{cousinweb}
\textbf{\gu{ઉપક્રમ}} (\emph{upakram}) is a tatsama, borrowed from Sanskrit \emph{upakrama}, \textquotedblleft{}approach, beginning, undertaking\textquotedblright{} (Monier-Williams).
\end{cousinweb}

\begin{grammarlens}[title={What you've built}]
One more word for saying what happens if.
\end{grammarlens}

\subsection*{Your turn}
\begin{itemize}
\raggedright
  \item \emph{Say it:} \emph{upakram}
  \item \emph{Say it:} \emph{upakram}, once more
  \item \emph{Say it:} \emph{jo ā upakram saphaḷ thaśe to bījā̃ gāmo paṇ te apnāvśe.} — If this initiative succeeds, other villages will take it up too.
  \item \emph{Recall:} say the Gujarati for to deteriorate, then the Gujarati for an outline, then say \emph{upakram} again
\end{itemize}

\begin{tcolorbox}[breakable,colback=teal!4,colframe=teal!35!black,title={Before you move on}]
What does \gu{ઉપક્રમ} mean? (\textquotedblleft{}An initiative\textquotedblright{}.) Say it once more.
\end{tcolorbox}
\section[\texorpdfstring{\gu{અભિયાન} (abhiyān)}{અભિયાન (abhiyān)}]{\gu{અભિયાન} (abhiyān) — a public drive; a campaign}

\label{lesson:GU-C388-abhiyan}



\begin{quote}
Before the new one: say the Gujarati for an outline, then the Gujarati for an initiative.
\end{quote}

\subsection*{You'll want to know: \gu{અભિયાન}}
\textbf{\gu{અભિયાન}} — \emph{abhiyān} — \textquotedblleft{}a public drive; a campaign\textquotedblright{}.

\begin{quote}
\gu{જો} \gu{આ} \gu{અભિયાન} \gu{ચાલશે} \gu{તો} \gu{ગામ} \gu{ચોખ્ખું} \gu{રહેશે}.
\end{quote}

\begin{quote}
\emph{jo ā abhiyān cālśe to gām cokhkhũ raheśe.}
\end{quote}

\begin{quote}
If this campaign keeps going, the village will stay clean.
\end{quote}

A neuter noun.

\begin{cousinweb}
\textbf{\gu{અભિયાન}} (\emph{abhiyān}) is borrowed from Sanskrit \emph{abhiyāna}, \textquotedblleft{}coming near, approaching\textquotedblright{} (Monier-Williams); \textquotedblleft{}a campaign\textquotedblright{} is its modern sense.
\end{cousinweb}

\begin{grammarlens}[title={What you've built}]
One more word for saying what happens if.
\end{grammarlens}

\subsection*{Your turn}
\begin{itemize}
\raggedright
  \item \emph{Say it:} \emph{abhiyān}
  \item \emph{Say it:} \emph{abhiyān}, once more
  \item \emph{Say it:} \emph{jo ā abhiyān cālśe to gām cokhkhũ raheśe.} — If this campaign keeps going, the village will stay clean.
  \item \emph{Recall:} say the Gujarati for an outline, then the Gujarati for an initiative, then say \emph{abhiyān} again
\end{itemize}

\begin{tcolorbox}[breakable,colback=teal!4,colframe=teal!35!black,title={Before you move on}]
What does \gu{અભિયાન} mean? (\textquotedblleft{}A public drive; a campaign\textquotedblright{}.) Say it once more.
\end{tcolorbox}
\section[\texorpdfstring{\gu{આરંભવું} (ārambhvũ)}{આરંભવું (ārambhvũ)}]{\gu{આરંભવું} (ārambhvũ) — to commence}

\label{lesson:GU-C388-arambhvu}



\begin{quote}
Before the new one: say the Gujarati for an initiative, then the Gujarati for a public drive.
\end{quote}

\subsection*{You'll want to know: \gu{આરંભવું}}
\textbf{\gu{આરંભવું}} — \emph{ārambhvũ} — \textquotedblleft{}to commence\textquotedblright{}.

\begin{quote}
\gu{જો} \gu{બધા} \gu{સમયસર} \gu{આવશે} \gu{તો} \gu{આપણે} \gu{કાર્યક્રમ} \gu{આરંભીશું}.
\end{quote}

\begin{quote}
\emph{jo badhā samaysar āvśe to āpaṇe kāryakram ārambhīśũ.}
\end{quote}

\begin{quote}
If everyone arrives on time, we will begin the programme.
\end{quote}

Past, with \gu{તેણે} (he or she …): \gu{તેણે} \gu{નવું} \gu{કામ} \gu{આરંભ્યું}, \emph{teṇe navũ kām ārambhyũ}.

\begin{cousinweb}
\textbf{\gu{આરંભવું}} (\emph{ārambhvũ}) is built on Sanskrit \emph{ārambha}, \textquotedblleft{}undertaking, beginning\textquotedblright{} (Monier-Williams), with the Gujarati verb ending.
\end{cousinweb}

\begin{grammarlens}[title={What you've built}]
One more word for saying what happens if.
\end{grammarlens}

\subsection*{Your turn}
\begin{itemize}
\raggedright
  \item \emph{Say it:} \emph{ārambhvũ}
  \item \emph{Say it:} \emph{ārambhvũ}, once more
  \item \emph{Say it:} \emph{jo badhā samaysar āvśe to āpaṇe kāryakram ārambhīśũ.} — If everyone arrives on time, we will begin the programme.
  \item \emph{Recall:} say the Gujarati for an initiative, then the Gujarati for a public drive, then say \emph{ārambhvũ} again
\end{itemize}

\begin{tcolorbox}[breakable,colback=teal!4,colframe=teal!35!black,title={Before you move on}]
What does \gu{આરંભવું} mean? (\textquotedblleft{}To commence\textquotedblright{}.) Say it once more.
\end{tcolorbox}
\section[\texorpdfstring{\gu{આદરવું} (ādarvũ)}{આદરવું (ādarvũ)}]{\gu{આદરવું} (ādarvũ) — to undertake}

\label{lesson:GU-C388-adarvu}



\begin{quote}
Before the new one: say the Gujarati for a public drive, then the Gujarati for to commence.
\end{quote}

\subsection*{You'll want to know: \gu{આદરવું}}
\textbf{\gu{આદરવું}} — \emph{ādarvũ} — \textquotedblleft{}to undertake\textquotedblright{}.

\begin{quote}
\gu{જો} \gu{તમે} \gu{આ} \gu{કામ} \gu{આદરશો} \gu{તો} \gu{હું} \gu{પૈસાની} \gu{મદદ} \gu{કરીશ}.
\end{quote}

\begin{quote}
\emph{jo tame ā kām ādarśo to hũ paisānī madad karīś.}
\end{quote}

\begin{quote}
If you take on this work, I will help with money.
\end{quote}

Past, with \gu{તેણે} (he or she …): \gu{તેણે} \gu{મોટું} \gu{કામ} \gu{આદર્યું}, \emph{teṇe moṭũ kām ādaryũ}.

\begin{cousinweb}
\textbf{\gu{આદરવું}} (\emph{ādarvũ}) is inherited from Sanskrit \emph{ādriyatē}, \textquotedblleft{}regards, is careful about\textquotedblright{}; CDIAL lists the Gujarati verb, \textquotedblleft{}to work with zeal, to begin\textquotedblright{} (CDIAL 1161). \textbf{\gu{આદર}} (\emph{ādar}), \textquotedblleft{}respect\textquotedblright{}, comes from the same root.
\end{cousinweb}

\begin{grammarlens}[title={What you've built}]
One more word for saying what happens if.
\end{grammarlens}

\subsection*{Your turn}
\begin{itemize}
\raggedright
  \item \emph{Say it:} \emph{ādarvũ}
  \item \emph{Say it:} \emph{ādarvũ}, once more
  \item \emph{Say it:} \emph{jo tame ā kām ādarśo to hũ paisānī madad karīś.} — If you take on this work, I will help with money.
  \item \emph{Recall:} say the Gujarati for a public drive, then the Gujarati for to commence, then say \emph{ādarvũ} again
\end{itemize}

\begin{tcolorbox}[breakable,colback=teal!4,colframe=teal!35!black,title={Before you move on}]
What does \gu{આદરવું} mean? (\textquotedblleft{}To undertake\textquotedblright{}.) Say it once more.
\end{tcolorbox}
